\documentclass[12pt]{article}

\usepackage{sbc-template}
\usepackage{graphicx,url}
\usepackage[utf8]{inputenc}
\usepackage[brazil]{babel}
\usepackage{amsmath}
\usepackage{amssymb}
\usepackage{booktabs}
\usepackage{multirow}

\sloppy

\title{EcoLex RAG: Sistema HyPA-RAG com Legal-BERTimbau para\\ Suporte Explicável a Decisão em Gestão Ambiental Brasileira}

\author{Delvek da Silva Venceslau de Sousa\inst{1}}

\address{Universidade Federal do Pará (UFPA) -- Belém, PA -- Brasil
  \email{delvek.sousa@icen.ufpa.br}
}

\begin{document}

\maketitle

\begin{abstract}
Brazilian environmental management demands frequent consultation of an extensive and hierarchically complex legislative corpus, a task that poses significant challenges to public administrators and environmental analysts. This paper presents EcoLex RAG, a Hybrid Parameter-Adaptive Retrieval-Augmented Generation (HyPA-RAG) system for explainable decision support in environmental management. The proposed architecture integrates three retrieval mechanisms -- BM25 (sparse), Legal-BERTimbau/FAISS (dense), and a legal Knowledge Graph -- unified via Reciprocal Rank Fusion, with adaptive parameters controlled by a query complexity classifier. Response generation is performed by Mistral 7B Instruct in FP16, with structured prompts requiring exact citation of articles, paragraphs, and clauses. An ablation study with 100 queries across 6 legislative acts demonstrated that the HyPA-RAG configuration achieves 65\% Law Retrieval and 48\% Citation Article Rate, with superior performance on medium-complexity queries (54\% CAR) but identified degradation on simple queries (40\% CAR vs. 50\% for BM25-only) using on average only 6.8 sources compared to 10.0 in static configurations, evidencing superior context efficiency. Results indicate that parametric adaptation enables more efficient RAG systems for the Brazilian legal-environmental domain.

\textbf{Keywords:} Retrieval-Augmented Generation, Environmental Legislation, HyPA-RAG, Legal-BERTimbau, Explainable Artificial Intelligence.
\end{abstract}

\begin{resumo}
A gestão ambiental brasileira demanda consulta frequente a um corpus legislativo extenso e hierarquicamente complexo, tarefa que impõe desafios significativos a gestores públicos e analistas ambientais. Este trabalho apresenta o EcoLex RAG, um sistema de Geração Aumentada por Recuperação com Parâmetros Híbridos Adaptativos (HyPA-RAG) voltado ao suporte explicável a decisão em gestão ambiental. A arquitetura proposta integra três mecanismos de recuperação -- BM25 (\textit{sparse}), Legal-BERTimbau/FAISS (\textit{dense}) e Knowledge Graph legal -- unificados via Reciprocal Rank Fusion, com parâmetros adaptativos controlados por um classificador de complexidade de consultas. A geração de respostas é realizada pelo Mistral 7B Instruct em FP16, com prompts estruturados que exigem citação exata de artigos, parágrafos e incisos. Um estudo de ablação com 100 consultas sobre 6 legislações demonstrou que a configuração HyPA-RAG alcança 65\% em Lei Retrieval e 48\% em Citation Article Rate, com desempenho superior em consultas de complexidade média (54\% CAR) mas com degradação identificada em consultas simples (40\% CAR vs. 50\% do BM25-only) utilizando em média apenas 6,8 fontes, contra 10,0 das configurações estáticas, evidenciando eficiência superior na economia de contexto. Os resultados indicam que a adaptação paramétrica viabiliza sistemas RAG mais eficientes para o domínio jurídico-ambiental brasileiro.

\textbf{Palavras-chave:} Geração Aumentada por Recuperação, Legislação Ambiental, HyPA-RAG, Legal-BERTimbau, Inteligência Artificial Explicável.
\end{resumo}

\section{Introdução}

O Brasil detém um dos arcabouços normativos ambientais mais extensos e sofisticados do planeta, abrangendo desde o Código Florestal (Lei 12.651/2012), que disciplina a proteção de Áreas de Preservação Permanente (APP) e Reservas Legais, até a Política Nacional do Meio Ambiente (Lei 6.938/1981), que institui o Sistema Nacional do Meio Ambiente (SISNAMA), passando pelo Sistema Nacional de Unidades de Conservação (Lei 9.985/2000), a Lei de Crimes Ambientais (Lei 9.605/1998) e diversas resoluções do Conselho Nacional do Meio Ambiente (CONAMA). A estrutura hierárquica dessas normas -- organizadas em artigos, parágrafos, incisos, alíneas e dispositivos remissivos entre leis distintas -- impõe uma carga cognitiva substancial a gestores públicos, analistas ambientais e operadores do direito ambiental que necessitam fundamentar suas decisões em dispositivos legais específicos \cite{robertson:2009}.

Esse cenário adquire relevância estratégica inédita no contexto pós-COP30 (Conferência das Nações Unidas sobre Mudanças Climáticas), realizada em Belém do Pará em novembro de 2025, evento que reposicionou o Brasil -- e particularmente a região amazônica -- no epicentro das discussões globais sobre governança climática e bioeconomia. As metas assumidas pelo país durante a COP30 exigem, em termos práticos, que órgãos ambientais estaduais e municipais operem com maior agilidade e precisão no enquadramento legal de atividades econômicas, licenciamentos e fiscalizações. A ausência de ferramentas computacionais adequadas para consulta inteligente a legislação constitui, portanto, um gargalo operacional com implicações diretas sobre a efetividade da governança ambiental brasileira.

No campo da Inteligência Artificial, a técnica de Geração Aumentada por Recuperação (RAG -- \textit{Retrieval-Augmented Generation}) \cite{lewis:2020} emergiu como paradigma dominante para a construção de sistemas de pergunta-resposta fundamentados em bases de conhecimento, mitigando o problema de alucinação factual em Modelos de Linguagem de Grande Escala (LLMs). Entretanto, a aplicação direta de arquiteturas RAG convencionais ao domínio jurídico-ambiental brasileiro enfrenta três lacunas fundamentais: (i) modelos de \textit{embedding} genéricos não capturam adequadamente a semântica de termos legais em português, como ``Área de Preservação Permanente'', ``Reserva Legal'' ou ``outorga de recurso hídrico''; (ii) a ausência de mecanismos que explorem a estrutura hierárquica intrínseca da legislação (Lei $>$ Capítulo $>$ Seção $>$ Artigo $>$ Parágrafo $>$ Inciso $>$ Alínea) desperdiça informação relacional valiosa; e (iii) a rigidez paramétrica de sistemas RAG tradicionais -- que operam com valores fixos de documentos recuperados independentemente da complexidade da consulta -- conduz a ineficiência tanto na economia de contexto quanto na qualidade da resposta.

Recentemente, \cite{kalra:2024} propuseram o HyPA-RAG (\textit{Hybrid Parameter-Adaptive Retrieval-Augmented Generation}), uma arquitetura que combina recuperação híbrida (\textit{sparse} + \textit{dense} + \textit{knowledge graph}) com adaptação paramétrica baseada na complexidade da consulta, demonstrando resultados promissores no domínio jurídico anglofono. Contudo, a transferência direta dessa arquitetura ao contexto legal brasileiro permanece inexplorada na literatura.

O presente trabalho propõe o \textbf{EcoLex RAG}, um sistema HyPA-RAG adaptado ao domínio jurídico-ambiental brasileiro, cujas contribuições principais são:

\begin{enumerate}
\item \textbf{Recuperação híbrida tripla adaptada ao domínio legal brasileiro}, integrando BM25 (\textit{sparse}), Legal-BERTimbau/FAISS (\textit{dense}) \cite{rufimelo:2022} e um Knowledge Graph legal com triplets extraídos da estrutura normativa, unificados via Reciprocal Rank Fusion (RRF) \cite{cormack:2009};
\item \textbf{Classificador de complexidade de consultas} baseado em heurísticas jurídico-linguísticas, que governa a adaptação de parâmetros ($top\_k$ e \textit{query rewrites}) por nível de complexidade (simples, médio, complexo);
\item \textbf{Pipeline de geração explicável} com Mistral 7B Instruct \cite{jiang:2023} em FP16, estruturado para produzir respostas com citação exata de dispositivos legais (artigo, parágrafo, inciso);
\item \textbf{Estudo de ablação abrangente} com 100 consultas distribuídas em três níveis de complexidade sobre 6 legislações ambientais, avaliando sistematicamente a contribuição de cada componente da arquitetura.
\end{enumerate}

O restante deste artigo está organizado da seguinte forma: a Seção 2 revisa os trabalhos relacionados; a Seção 3 detalha a metodologia e a arquitetura do sistema; a Seção 4 descreve o protocolo experimental; a Seção 5 apresenta e discute os resultados; e a Seção 6 conclui o trabalho com direcionamentos futuros.

\section{Trabalhos Relacionados}

Esta seção situa o EcoLex RAG no panorama da literatura em três eixos: Geração Aumentada por Recuperação, modelos de linguagem para o português jurídico e aplicações de IA na gestão ambiental.

\subsection{Geração Aumentada por Recuperação (RAG)}

\cite{lewis:2020} formalizaram o paradigma RAG ao demonstrar que a concatenação de documentos recuperados de uma base de conhecimento ao \textit{prompt} de um modelo \textit{seq2seq} (\textit{sequence-to-sequence}) produz respostas mais factuais do que modelos puramente generativos. Desde então, a arquitetura RAG tornou-se o padrão \textit{de facto} para sistemas de pergunta-resposta em domínios especializados, sendo adotada em contextos médicos, jurídicos e financeiros.

A recuperação \textit{sparse} via BM25 \cite{robertson:2009} permanece como \textit{baseline} robusto, particularmente eficaz na correspondência de termos exatos -- propriedade crítica no domínio legal, onde identificadores normativos como ``Art. 4\textsuperscript{o}'' ou ``Resolução CONAMA 357'' constituem tokens de alta especificidade. Paralelamente, a recuperação \textit{dense} baseada em modelos de \textit{embedding} captura relações semânticas que escapam a correspondência léxica, sendo especialmente útil em consultas formuladas em linguagem natural por usuários não especializados.

\subsection{HyPA-RAG: Recuperação Híbrida com Parâmetros Adaptativos}

\cite{kalra:2024} introduziram o HyPA-RAG (\textit{Hybrid Parameter-Adaptive Retrieval-Augmented Generation}), uma extensão do paradigma RAG que endereça duas limitações fundamentais das arquiteturas convencionais. Primeiramente, o HyPA-RAG combina três modalidades de recuperação -- \textit{sparse} (BM25), \textit{dense} (\textit{embeddings}) e \textit{graph-based} (Knowledge Graph) -- mediante Reciprocal Rank Fusion (RRF), proposta originalmente por \cite{cormack:2009}. Em segundo lugar, a arquitetura classifica cada consulta em níveis de complexidade e adapta dinamicamente os parâmetros de recuperação (número de documentos, reformulações de consulta), evitando tanto a subrecuperação em consultas complexas quanto o desperdício de contexto em consultas simples. Os autores demonstraram ganhos significativos em benchmarks jurídicos em língua inglesa, particularmente em consultas que exigem raciocínio \textit{multi-hop} entre dispositivos normativos.

A fórmula do RRF é formalizada na Equação~\ref{eq:rrf_detail} (Seção 3.4). A fusão por RRF dispensa calibração de pesos entre modelos heterogêneos, propriedade particularmente vantajosa quando os escores dos diferentes \textit{retrievers} residem em escalas incomparáveis, uma vez que utiliza apenas as posições no ranking, com constante de suavização $k = 60$ conforme o \textit{paper} original \cite{cormack:2009}.

\subsection{Modelos de Linguagem para o Português Jurídico}

O BERTimbau \cite{souza:2020} estabeleceu o primeiro modelo BERT pré-treinado especificamente para o português brasileiro, treinado sobre um corpus de 2,7 bilhões de tokens extraídos do \textit{BrWaC} (\textit{Brazilian Web as Corpus}). Subsequentemente, \cite{rufimelo:2022} apresentaram o Legal-BERTimbau, uma variante do BERTimbau submetida a treinamento contínuo sobre um corpus jurídico em português, demonstrando ganhos consistentes em tarefas de \textit{Semantic Textual Similarity} (STS) no domínio legal. O modelo \texttt{rufimelo/Legal-BERTimbau-sts-base-ma}, com dimensão de \textit{embedding} de 768, é particularmente adequado para a geração de representações vetoriais de trechos legislativos brasileiros, capturando nuances semânticas que modelos genéricos não discriminam.

Para a geração de respostas, o Mistral 7B \cite{jiang:2023} emergiu como uma alternativa eficiente a modelos de maior escala, oferecendo desempenho competitivo com LLMs de 13 bilhões de parâmetros graças ao mecanismo de atenção de janela deslizante (\textit{Sliding Window Attention}) e agrupamento de consultas (\textit{Grouped-Query Attention}). A variante Instruct, ajustada via \textit{Reinforcement Learning from Human Feedback} (RLHF), demonstra capacidade superior de seguir instruções estruturadas, propriedade essencial para a geração de respostas com citação legal obrigatória.

\subsection{IA Aplicada a Gestão Ambiental e Explicabilidade}

A intersecção entre IA e gestão ambiental tem avançado substancialmente, com aplicações que vão desde a detecção de desmatamento por visão computacional até a modelagem preditiva de dinâmicas climáticas \cite{wcama:2025}. Contudo, a vertente de suporte a decisão jurídico-ambiental permanece subexplorada. No contexto da governança pública, a demanda por Inteligência Artificial Explicável (XAI -- \textit{Explainable Artificial Intelligence}) tornou-se imperativa: decisões administrativas fundamentadas em recomendações de IA devem ser auditáveis e rastreáveis aos dispositivos normativos que as embasam. Sistemas RAG com citação de fontes representam, nesse sentido, uma forma nativa de explicabilidade, na medida em que cada afirmação da resposta pode ser verificada contra o trecho legislativo original.

A edição de 2025 do WCAMA evidenciou a maturidade da ecoinformática brasileira e a crescente adoção de técnicas como RAG e XAI para problemas ambientais, sinalizando que o campo está receptivo a contribuições que integrem processamento de linguagem natural, representação de conhecimento legal e suporte a decisão governamental.

\section{Metodologia}

Esta seção detalha a arquitetura do sistema EcoLex RAG, abordando cada componente do pipeline HyPA-RAG: classificação de complexidade, recuperação tripla, fusão de rankings, adaptação paramétrica e geração de respostas explicáveis.

\subsection{Visão Geral da Arquitetura}

O EcoLex RAG implementa um pipeline HyPA-RAG completo, cujo fluxo de processamento é ilustrado na Figura~\ref{fig:arquitetura}. Uma consulta em linguagem natural ingressada pelo gestor ambiental percorre cinco estágios sequenciais: (1) classificação de complexidade; (2) adaptação de parâmetros; (3) recuperação tripla com fusão; (4) geração de resposta com LLM; e (5) entrega de resposta explicável com fontes exatas.

\begin{figure}[ht]
\centering
\begin{small}
\begin{verbatim}
  Consulta do Gestor Ambiental
              |
              v
  +----------------------------+
  | Classificador de           |
  | Complexidade               |
  | (heurísticas linguísticas) |
  +----------------------------+
              |
              v
  +----------------------------+
  | Parâmetros Adaptativos     |
  | k=5/10/15                  |
  | rewrites=0/1/3             |
  +----------------------------+
              |
    +---------+---------+
    |         |         |
    v         v         v
+--------+ +--------+ +--------+
|  BM25  | | Legal- | | Know-  |
| Okapi  | | BERT-  | | ledge  |
|(sparse)| | imbau  | | Graph  |
|        | | +FAISS | |(triplets|
|        | |(dense) | | legais)|
+---+----+ +---+----+ +---+----+
    |         |         |
    +---------+---------+
              |
              v
  +----------------------------+
  | Reciprocal Rank Fusion     |
  | RRF(d) = Sum 1/(k+rank(d))|
  +----------------------------+
              |
              v
  +----------------------------+
  | Mistral 7B Instruct (FP16) |
  | Prompt: citação obrigatória|
  | Estrutura: FUNDAMENTACAO   |
  |   LEGAL > ANALISE >        |
  |   CONCLUSAO                |
  +----------------------------+
              |
              v
  +----------------------------+
  | Resposta Explicável        |
  | + Fontes com Lei/Art/Par   |
  | + Metadados HyPA           |
  +----------------------------+
\end{verbatim}
\end{small}
\caption{Arquitetura do pipeline HyPA-RAG do EcoLex RAG. A consulta é classificada em três níveis de complexidade, que governam os parâmetros adaptativos dos três \textit{retrievers} operando em paralelo. A fusão via RRF alimenta o LLM com contexto legislativo ranqueado, e a resposta é entregue com citação obrigatória de fontes.}
\label{fig:arquitetura}
\end{figure}

\subsection{Classificador de Complexidade}

O classificador de complexidade constitui o componente que diferencia a abordagem HyPA-RAG de sistemas RAG tradicionais com parâmetros estáticos. Dado que consultas jurídicas variam enormemente em profundidade -- de perguntas factuais diretas (``O que é APP?'') a consultas comparativas \textit{multi-hop} (``Compare as exigências de APP e Reserva Legal para propriedade na Amazônia Legal'') --, a adaptação paramétrica permite alocar recursos computacionais proporcionalmente a demanda informacional de cada consulta.

O classificador opera sobre heurísticas linguísticas específicas do domínio jurídico-ambiental, computando um \textit{complexity score} $s$ baseado em seis indicadores extraídos da consulta $q$:

Seja $q$ uma consulta tokenizada em palavras $w_1, w_2, \ldots, w_n$. O escore de complexidade é calculado como:

\begin{equation}
s(q) = \min(\lambda_L, 3) + 2 \cdot \mathbb{1}_{C}(q) + \mathbb{1}_{D}(q) + 2 \cdot \mathbb{1}_{M}(q) + \mathbb{1}_{?}(q) + \mathbb{1}_{\wedge}(q) + \mathbb{1}_{|q|>30}
\label{eq:complexity}
\end{equation}

\noindent onde:
\begin{itemize}
\item $\lambda_L = |\{t \in \mathcal{T}_L : t \subset q\}|$ é a contagem de termos legais reconhecidos, limitada a 3, com $\mathcal{T}_L$ sendo o léxico de 36 termos jurídico-ambientais (e.g., ``reserva legal'', ``licenciamento'', ``CONAMA'', ``outorga'', ``recurso hídrico'');
\item $\mathbb{1}_{C}(q)$ indica presença de marcadores comparativos (``compare'', ``diferença'', ``versus'');
\item $\mathbb{1}_{D}(q)$ indica presença de marcadores condicionais (``se'', ``caso'', ``na hipótese de'');
\item $\mathbb{1}_{M}(q)$ indica presença de marcadores \textit{multi-hop} (``adicionalmente'', ``cumulativamente'', ``à luz de'');
\item $\mathbb{1}_{?}(q)$ indica mais de um ponto de interrogação;
\item $\mathbb{1}_{\wedge}(q)$ indica duas ou mais conjunções (``e'', ``ou'', ``mas'', ``porém'');
\item $\mathbb{1}_{|q|>30}$ indica que a consulta excede 30 palavras.
\end{itemize}

A classificação final mapeia o escore a três níveis:

\begin{equation}
\text{complexidade}(q) = \begin{cases} \texttt{simple} & \text{se } s(q) \leq 2 \\ \texttt{medium} & \text{se } 3 \leq s(q) \leq 5 \\ \texttt{complex} & \text{se } s(q) \geq 6 \end{cases}
\label{eq:classification}
\end{equation}

\subsection{Recuperação Tripla}

O módulo de recuperação opera simultaneamente três estratégias complementares para maximizar a cobertura e a precisão na identificação de trechos legislativos relevantes.

\textbf{BM25 Okapi (\textit{sparse}).} O algoritmo BM25 \cite{robertson:2009} realiza recuperação baseada em correspondência léxica, computando a relevância de cada \textit{chunk} legislativo em função da frequência dos termos da consulta, ponderada pela frequência inversa no corpus e pelo comprimento do documento. A tokenização é realizada por \textit{whitespace splitting} em \textit{lowercase}, sem \textit{stemming}, preservando a integridade de identificadores normativos (e.g., ``Art. 4\textsuperscript{o}'', ``CONAMA 357''). O índice BM25 é construído sobre 473 \textit{chunks} extraídos de 6 legislações, armazenados em formato JSON.

\textbf{Legal-BERTimbau + FAISS (\textit{dense}).} A recuperação densa utiliza o modelo \texttt{rufimelo/Legal-BERTimbau-sts-base-ma} \cite{rufimelo:2022} para gerar \textit{embeddings} de dimensão $d = 768$ para cada \textit{chunk} legislativo. Os vetores são normalizados via $L_2$ e indexados em uma estrutura FAISS \texttt{IndexFlatIP} (\textit{Inner Product}), que computa a similaridade por produto interno (equivalente a similaridade cosseno após normalização). A busca por similaridade retorna os $k$ \textit{chunks} mais próximos ao vetor da consulta no espaço semântico.

\textbf{Knowledge Graph Legal (\textit{graph-based}).} O Knowledge Graph (KG) é construído a partir da estrutura hierárquica intrínseca da legislação brasileira. Um extrator baseado em expressões regulares processa o texto legislativo e gera triplets no formato $(sujeito, relação, objeto)$, onde as relações capturam cinco tipos de vínculos:
\begin{itemize}
\item \texttt{pertence\_a}: artigo pertence a uma lei;
\item \texttt{detalha}: parágrafo detalha um artigo;
\item \texttt{define}: inciso define um conceito dentro de um artigo;
\item \texttt{definido\_em}: conceito legal é definido em um artigo específico;
\item \texttt{estabelece\_medida}: artigo estabelece medida quantitativa (metros, hectares, percentuais).
\end{itemize}

A busca no KG opera por correspondência de termos entre a consulta e os campos dos triplets, com \textit{boost} de 0,3 para \textit{matches} no campo \texttt{subject} e 0,2 para \textit{matches} no campo \texttt{object}. Cada triplet retornado é formatado como \textit{chunk} textual contendo a relação e o contexto associado (limitado a 300 caracteres).

\subsection{Reciprocal Rank Fusion}

Os resultados dos três \textit{retrievers} são unificados via Reciprocal Rank Fusion (RRF) \cite{cormack:2009}. Para cada documento $d$ presente em ao menos uma lista de resultados, o escore RRF é calculado como:

\begin{equation}
\text{RRF}_{\text{score}}(d) = \sum_{r \in \{BM25, Dense, KG\}} \frac{1}{k + \text{rank}_r(d)}
\label{eq:rrf_detail}
\end{equation}

\noindent com $k = 60$ (constante de suavização padrão). Documentos ausentes em uma lista não contribuem para a soma. Os documentos fusionados são ordenados por escore RRF decrescente, e os top-$n$ são selecionados como contexto para o LLM. A deduplicação é realizada por identidade dos primeiros 200 caracteres do conteúdo, garantindo que \textit{chunks} idênticos recuperados por \textit{retrievers} distintos não sejam contados duplamente. O metadado \texttt{retrieval\_method} do \textit{chunk} fusionado registra todos os métodos que o recuperaram (e.g., ``bm25+semantic''), viabilizando auditoria do pipeline.

\subsection{Parâmetros Adaptativos}

A adaptação paramétrica constitui o núcleo diferenciador da arquitetura HyPA-RAG em relação a sistemas RAG convencionais. Com base na classificação de complexidade da consulta (Seção 3.2), o pipeline ajusta dois parâmetros: (i) o número de documentos recuperados por cada \textit{retriever} ($top\_k$) e (ii) o número de reformulações de consulta (\textit{query rewrites}) geradas pelo LLM para ampliar a cobertura léxica e semântica da busca. A Tabela~\ref{tab:parâmetros} apresenta a configuração adotada.

\begin{table}[ht]
\centering
\caption{Parâmetros adaptativos do HyPA-RAG por nível de complexidade.}
\label{tab:parâmetros}
\begin{tabular}{cccp{5.5cm}}
\toprule
\textbf{Complexidade} & \textbf{$top\_k$} & \textbf{\textit{Query Rewrites}} & \textbf{Exemplo de consulta} \\
\midrule
Simples & 5 & 0 & ``O que é APP?'' \\
Média & 10 & 1 & ``Qual a faixa de APP para rio de 10m?'' \\
Complexa & 15 & 3 & ``Compare as exigências de APP e Reserva Legal para propriedade na Amazônia Legal'' \\
\bottomrule
\end{tabular}
\end{table}

A redução do $top\_k$ para consultas simples economiza contexto no \textit{prompt} do LLM, diminuindo o custo computacional de inferência e mitigando o risco de ``diluição'' da resposta por \textit{chunks} irrelevantes. As reformulações de consulta (\textit{query rewrites}) são geradas pelo próprio Mistral 7B com temperatura de 0,3, produzindo variações linguísticas da pergunta original (e.g., sinônimos de termos legais, reestruturações sintáticas) que ampliam o espaço de busca nos \textit{retrievers}.

\subsection{Geração com Mistral 7B Instruct}

A geração de respostas é realizada pelo modelo Mistral 7B Instruct v0.3 \cite{jiang:2023}, carregado em precisão FP16 (\textit{float16}) com mapeamento automático de dispositivo (\texttt{device\_map="auto"}), exigindo aproximadamente 14 GB de VRAM. O pipeline de geração utiliza a biblioteca Hugging Face \texttt{transformers} com os seguintes hiperparâmetros: \texttt{max\_new\_tokens=1024}, \texttt{temperature=0.1} e \texttt{do\_sample=True}.

O \textit{prompt} do sistema é estruturado com seis regras obrigatórias que constituem o mecanismo de explicabilidade do EcoLex RAG: (1) responder exclusivamente com base nos trechos legislativos fornecidos; (2) citar a fonte exata (lei, artigo, parágrafo, inciso); (3) declarar explicitamente a insuficiência de informação quando aplicável; (4) utilizar linguagem técnica acessível a gestores públicos; (5) estruturar a resposta em FUNDAMENTAÇÃO LEGAL, ANÁLISE e CONCLUSÃO; e (6) não inventar dispositivos legais ausentes do contexto.

O contexto legislativo é formatado como sequência de \textit{chunks}, cada um prefixado com os metadados \texttt{[lei | artigo]}, limitados a 800 caracteres e até 3.000 tokens estimados no total, garantindo operação dentro da janela de contexto do modelo sem exceder a capacidade da VRAM.

\subsection{Knowledge Graph Legal: Extração de Triplets}

A construção do Knowledge Graph legal é realizada por um extrator baseado em expressões regulares que processa o texto integral de cada legislação. O extrator reconhece cinco padrões estruturais:

\begin{enumerate}
\item \textbf{Artigos}: padrão \texttt{Art.\textbackslash s*\textbackslash d+[\textbackslash w-]*} identifica todos os artigos e associa cada um a lei correspondente via triplet $(Art.\ X, pertence\_a, Lei\ Y)$;
\item \textbf{Parágrafos}: padrão \texttt{\S\textbackslash s*\textbackslash d+} identifica parágrafos e os vincula ao artigo pai via $(Art.\ X, \S Y, detalha, Art.\ X)$;
\item \textbf{Incisos}: identificadores romanos (I, II, III, \ldots) são extraídos e vinculados ao artigo via $(Art.\ X, Inciso\ Y, define, Art.\ X)$;
\item \textbf{Definições}: padrões como ``entende-se por'', ``considera-se'' e ``define-se como'' extraem conceitos definidos na legislação;
\item \textbf{Medidas quantitativas}: padrões numéricos seguidos de unidades (metros, hectares, percentuais) extraem limites quantitativos estabelecidos pela norma.
\end{enumerate}

Sobre as 6 legislações da base de conhecimento, o extrator gerou 1.714 triplets, constituindo uma representação estruturada das relações normativas que complementa a recuperação textual pura.

\section{Experimentos}

Esta seção descreve a base de conhecimento, o \textit{dataset} de avaliação, as configurações de ablação, as métricas empregadas e o ambiente de \textit{hardware}.

\subsection{Base de Conhecimento}

A base de conhecimento do EcoLex RAG compreende 6 legislações ambientais brasileiras, segmentadas em 473 \textit{chunks} textuais e representadas por 1.714 triplets no Knowledge Graph. A Tabela~\ref{tab:legislações} detalha a composição da base.

\begin{table}[ht]
\centering
\caption{Legislações que compõem a base de conhecimento do EcoLex RAG.}
\label{tab:legislações}
\begin{tabular}{lcp{5cm}}
\toprule
\textbf{Legislação} & \textbf{Tipo} & \textbf{Escopo Temático} \\
\midrule
Código Florestal (Lei 12.651/2012) & Lei Federal & APP, Reserva Legal, supressão de vegetação \\
PNMA (Lei 6.938/1981) & Lei Federal & Política Nacional do Meio Ambiente, SISNAMA, licenciamento \\
SNUC (Lei 9.985/2000) & Lei Federal & Unidades de Conservação, categorias de manejo \\
Lei de Crimes Ambientais (Lei 9.605/1998) & Lei Federal & Sanções penais e administrativas ambientais \\
Resolução CONAMA 357/2005 & Resolução & Classificação de corpos d'água, padrões de qualidade \\
Resolução CONAMA 430/2011 & Resolução & Condições e padrões de lançamento de efluentes \\
\bottomrule
\end{tabular}
\end{table}

\subsection{Dataset de Avaliação}

O \textit{dataset} de avaliação foi construído especificamente para este estudo, contendo 100 consultas em linguagem natural sobre as 6 legislações da base. Cada entrada do \textit{dataset} contém: (i) a pergunta em linguagem natural; (ii) a resposta de referência com fundamentação legal; (iii) a lei esperada nos resultados de recuperação; (iv) o artigo esperado; (v) o nível de complexidade; e (vi) a categoria temática. A distribuição por complexidade e legislação é apresentada na Tabela~\ref{tab:dataset}.

\begin{table}[ht]
\centering
\caption{Distribuição do \textit{dataset} de avaliação por complexidade e legislação.}
\label{tab:dataset}
\begin{tabular}{lcccc}
\toprule
\textbf{Legislação} & \textbf{Simples} & \textbf{Média} & \textbf{Complexa} & \textbf{Total} \\
\midrule
Código Florestal & 6 & 12 & 12 & 30 \\
PNMA & 6 & 5 & 5 & 16 \\
SNUC & 6 & 5 & 5 & 16 \\
Crimes Ambientais & 10 & 3 & 3 & 16 \\
CONAMA 357/2005 & 6 & 6 & 0 & 12 \\
CONAMA 430/2011 & 6 & 4 & 0 & 10 \\
\midrule
\textbf{Total} & \textbf{40} & \textbf{35} & \textbf{25} & \textbf{100} \\
\bottomrule
\end{tabular}
\end{table}

A preponderância do Código Florestal (30 perguntas) reflete sua centralidade na prática da gestão ambiental e sua maior complexidade estrutural (artigos com múltiplos parágrafos, incisos e alíneas). As consultas complexas concentram-se nas leis de maior inter-relação normativa (Código Florestal, PNMA, SNUC, Crimes Ambientais), enquanto as resoluções CONAMA, de escopo mais restrito, são cobertas por consultas simples e médias.

\subsection{Configurações de Ablação}

O estudo de ablação avalia 6 configurações que isolam e combinam progressivamente os componentes do pipeline, permitindo mensurar a contribuição individual e sinérgica de cada \textit{retriever}:

\begin{enumerate}
\item \textbf{BM25-only}: apenas recuperação \textit{sparse} via BM25 Okapi;
\item \textbf{Semantic-only}: apenas recuperação \textit{dense} via Legal-BERTimbau + FAISS;
\item \textbf{KG-only}: apenas recuperação via Knowledge Graph legal;
\item \textbf{BM25+Semantic}: fusão via RRF dos \textit{retrievers} BM25 e semântico (sem KG);
\item \textbf{Híbrido fixo}: fusão via RRF dos três \textit{retrievers} com parâmetros estáticos ($k = 10$, \textit{rewrites} $= 0$);
\item \textbf{HyPA-RAG}: fusão via RRF dos três \textit{retrievers} com parâmetros adaptativos (Tabela~\ref{tab:parâmetros}).
\end{enumerate}

As configurações 1--5 operam com $k = 10$ e 0 \textit{rewrites} em todas as consultas, enquanto a configuração 6 (HyPA-RAG) adapta estes parâmetros conforme a complexidade. Todas utilizam o mesmo LLM (Mistral 7B FP16) e o mesmo \textit{prompt} de geração.

\subsection{Métricas de Avaliação}

Quatro métricas são empregadas, abrangendo tanto a etapa de recuperação quanto a de geração:

\begin{itemize}
\item \textbf{Lei Retrieval Rate (LRR)}: fração de consultas cujos resultados de recuperação contêm ao menos um \textit{chunk} da lei esperada;
\item \textbf{Article Retrieval Rate (ARR)}: fração de consultas cujos resultados contêm o artigo específico esperado;
\item \textbf{Citation Law Rate (CLR)}: fração de respostas geradas que citam corretamente a lei esperada no corpo textual;
\item \textbf{Citation Article Rate (CAR)}: fração de respostas que citam o artigo específico esperado.
\end{itemize}

As duas primeiras métricas avaliam a qualidade do \textit{retriever}, enquanto as duas últimas avaliam a qualidade \textit{end-to-end} do pipeline (recuperação + geração). Adicionalmente, reporta-se o número médio de fontes (Fontes) e o tempo médio de resposta em segundos.

\subsection{Ambiente de Hardware}

Todos os experimentos foram conduzidos em uma estação de trabalho equipada com GPU NVIDIA GeForce RTX 5090 (24 GB VRAM), processador Intel Core Ultra 9 275HX e 64 GB de RAM DDR5. O Mistral 7B Instruct v0.3 foi carregado em precisão FP16, ocupando aproximadamente 14 GB de VRAM. A avaliação completa (600 execuções: 100 consultas $\times$ 6 configurações) consumiu aproximadamente 4,6 horas (16.610 segundos).

\section{Resultados e Discussão}

\subsection{Resultados Gerais}

A Tabela~\ref{tab:resultados_gerais} apresenta os resultados agregados das 6 configurações de ablação sobre as 100 consultas do \textit{dataset}.

\begin{table}[ht]
\centering
\caption{Resultados do estudo de ablação (métricas agregadas sobre 100 consultas).}
\label{tab:resultados_gerais}
\begin{tabular}{lcccccc}
\toprule
\textbf{Configuração} & \textbf{Lei Ret.} & \textbf{Art. Ret.} & \textbf{Cit. Lei} & \textbf{Cit. Art.} & \textbf{Fontes} & \textbf{Tempo(s)} \\
\midrule
BM25-only & 65\% & 70\% & 63\% & 45\% & 10,0 & 26,2 \\
Semantic-only & 61\% & 53\% & 61\% & 40\% & 10,0 & 24,6 \\
KG-only & 56\% & 35\% & 58\% & 31\% & 10,0 & 20,3 \\
BM25+Semantic & 65\% & 66\% & 65\% & 48\% & 10,0 & 27,6 \\
Híbrido fixo & 65\% & 59\% & 66\% & 49\% & 10,0 & 23,8 \\
\textbf{HyPA-RAG} & \textbf{65\%} & 56\% & 65\% & 48\% & \textbf{6,8} & 43,5 \\
\bottomrule
\end{tabular}
\end{table}

\subsection{Resultados por Complexidade: Consultas Complexas}

A Tabela~\ref{tab:complexas} apresenta o \textit{breakdown} de desempenho para as 25 consultas de complexidade alta, que envolvem raciocínio \textit{multi-hop}, comparações entre legislações e cenários condicionais.

\begin{table}[ht]
\centering
\caption{Desempenho nas 25 consultas complexas (\textit{multi-hop}, comparativas).}
\label{tab:complexas}
\begin{tabular}{lcccc}
\toprule
\textbf{Configuração} & \textbf{Lei Ret.} & \textbf{Art. Ret.} & \textbf{Cit. Lei} & \textbf{Cit. Art.} \\
\midrule
BM25-only & 68\% & 60\% & 68\% & 40\% \\
Semantic-only & 60\% & 48\% & 64\% & 44\% \\
KG-only & 60\% & 20\% & 64\% & 36\% \\
BM25+Semantic & 64\% & 52\% & 68\% & 52\% \\
Híbrido fixo & 64\% & 36\% & 72\% & 48\% \\
\textbf{HyPA-RAG} & 64\% & 52\% & 68\% & \textbf{52\%} \\
\bottomrule
\end{tabular}
\end{table}

\subsection{Resultados por Complexidade: Consultas Médias}

A Tabela~\ref{tab:médias} apresenta os resultados para as 35 consultas de complexidade média, tipicamente envolvendo condicionais e artigos com parágrafos.

\begin{table}[ht]
\centering
\caption{Desempenho nas 35 consultas de complexidade média (condicionais, parágrafos).}
\label{tab:médias}
\begin{tabular}{lcccc}
\toprule
\textbf{Configuração} & \textbf{Lei Ret.} & \textbf{Art. Ret.} & \textbf{Cit. Lei} & \textbf{Cit. Art.} \\
\midrule
BM25-only & 66\% & 63\% & 63\% & 43\% \\
Semantic-only & 60\% & 43\% & 60\% & 34\% \\
KG-only & 57\% & 31\% & 54\% & 23\% \\
BM25+Semantic & 66\% & 63\% & 66\% & 46\% \\
Híbrido fixo & 66\% & 63\% & 66\% & 49\% \\
\textbf{HyPA-RAG} & 66\% & 63\% & 63\% & \textbf{54\%} \\
\bottomrule
\end{tabular}
\end{table}

\subsection{Resultados por Complexidade: Consultas Simples}

A Tabela~\ref{tab:simples} apresenta os resultados para as 40 consultas de complexidade simples, que envolvem perguntas factuais diretas sobre um único dispositivo legal.

\begin{table}[ht]
\centering
\caption{Desempenho nas 40 consultas simples (factuais, dispositivo único).}
\label{tab:simples}
\begin{tabular}{lcccc}
\toprule
\textbf{Configuração} & \textbf{Lei Ret.} & \textbf{Art. Ret.} & \textbf{Cit. Lei} & \textbf{Cit. Art.} \\
\midrule
BM25-only & 62\% & 82\% & 60\% & 50\% \\
Semantic-only & 62\% & 65\% & 60\% & 42\% \\
KG-only & 52\% & 48\% & 57\% & 35\% \\
BM25+Semantic & 65\% & 78\% & 62\% & 48\% \\
Híbrido fixo & 65\% & 70\% & 62\% & 50\% \\
\textbf{HyPA-RAG} & 65\% & 52\% & 65\% & 40\% \\
\bottomrule
\end{tabular}
\end{table}

\subsection{Exemplo Qualitativo de Resposta}

Para ilustrar a explicabilidade do sistema, apresenta-se um exemplo de resposta gerada pelo HyPA-RAG para uma consulta de complexidade média:

\textbf{Consulta:} ``Qual a largura mínima da faixa de APP para cursos d'água naturais com menos de 10 metros de largura?''

\textbf{Complexidade classificada:} Média ($s = 3$, $top\_k = 10$, 1 \textit{rewrite})

\textbf{Resposta gerada (excerto):}

\begin{quote}
\textit{FUNDAMENTAÇÃO LEGAL: Conforme o Art. 4º, inciso I, alínea `a' do Código Florestal (Lei 12.651/2012), as faixas marginais de qualquer curso d'água natural perene e intermitente, excluídos os efêmeros, desde a borda da calha do leito regular, devem ter largura mínima de 30 (trinta) metros para cursos d'água de menos de 10 (dez) metros de largura.}

\textit{ANÁLISE: O dispositivo legal é claro ao estabelecer a faixa de proteção de APP com base na largura do curso d'água, utilizando critério objetivo mensurável.}

\textit{CONCLUSÃO: A faixa mínima de APP é de 30 metros, conforme Art. 4º, I, `a' da Lei 12.651/2012.}
\end{quote}

\textbf{Fontes recuperadas:} 3 \textit{chunks} do Código Florestal (Art. 4º), via BM25+Semantic. A resposta cita corretamente a lei, o artigo, o inciso e a alínea, demonstrando a rastreabilidade exigida pelo \textit{prompt} de explicabilidade.

\subsection{Discussão}

Os resultados do estudo de ablação permitem extrair cinco constatações fundamentais sobre a dinâmica dos componentes do EcoLex RAG.

\textbf{Superioridade do BM25 para termos legais exatos.} A configuração BM25-only alcança o melhor Article Retrieval Rate global (70\%), superando em larga margem as abordagens Semantic-only (53\%) e KG-only (35\%). Este resultado é esperado e corrobora a literatura \cite{robertson:2009}: identificadores normativos como ``Art. 4\textsuperscript{o}'' e ``CONAMA 357'' são tokens de alta especificidade cuja correspondência léxica exata é mais eficaz do que a proximidade no espaço de \textit{embeddings}. A recuperação \textit{dense}, por sua vez, captura relações semânticas úteis mas dilui a precisão ao trazer \textit{chunks} tematicamente próximos porém normativamente distintos.

\textbf{Knowledge Graph isolado e limitado, mas contribui na fusão.} O KG-only apresenta o pior desempenho em todas as métricas (56\% LRR, 35\% ARR, 31\% CAR), evidenciando que a recuperação puramente baseada em correspondência de termos nos triplets é insuficiente para cobrir a diversidade léxica das consultas. Contudo, sua contribuição na fusão tripla é perceptível: o Híbrido fixo (BM25+Semantic+KG) alcança Citation Law Rate de 66\%, superior aos 65\% do BM25+Semantic (sem KG), indicando que os triplets do KG adicionam informação estrutural complementar que melhora a capacidade do LLM de citar legislações corretamente.

\textbf{HyPA-RAG alcança desempenho competitivo com significativa economia de fontes.} Um aspecto relevante dos resultados reside na comparação entre o HyPA-RAG e as configurações estáticas em termos de eficiência. O HyPA-RAG atinge 65\% em Lei Retrieval e 48\% em Citation Article Rate -- valores competitivos com o Híbrido fixo (65\% e 49\%) e o BM25+Semantic (65\% e 48\%) -- utilizando em média apenas \textbf{6,8 fontes} contra 10,0 das demais configurações. Esta redução de 32\% no número de fontes traduz-se em: (i) menor consumo de tokens na janela de contexto do LLM; (ii) menor custo computacional de inferência; e (iii) menor risco de diluição da resposta por \textit{chunks} irrelevantes. A adaptação paramétrica permite que consultas simples ($k = 5$, 0 \textit{rewrites}) operem de forma leve, enquanto consultas complexas ($k = 15$, 3 \textit{rewrites}) recebem cobertura ampliada. Contudo, conforme discutido na análise das consultas simples (Tabela~\ref{tab:simples}), esta economia tem um custo mensurável em Article Retrieval Rate, configurando um \textit{tradeoff} que deve ser ponderado conforme o cenário de aplicação.

\textbf{Empate em consultas complexas, superioridade em consultas médias.} Nas 25 consultas complexas (Tabela~\ref{tab:complexas}), o HyPA-RAG e o BM25+Semantic empatam em Citation Article Rate (52\%), ambos superando o Híbrido fixo (48\%). Este resultado sugere que, para consultas que exigem raciocínio \textit{multi-hop}, a combinação de BM25 com recuperação semântica é o fator determinante, e a inclusão do KG na fusão não prejudica o desempenho quando os parâmetros são adaptados. Nas 35 consultas médias (Tabela~\ref{tab:médias}), o HyPA-RAG atinge \textbf{54\% em Citation Article Rate}, superando todas as demais configurações -- incluindo o Híbrido fixo (49\%) e o BM25+Semantic (46\%). A reformulação de consulta ($n = 1$ para complexidade média) e a calibração de $k = 10$ emergem como a combinação mais eficaz para este estrato de dificuldade, indicando que a adaptação paramétrica agrega maior valor precisamente no espectro intermediário de complexidade, onde nem os parâmetros mínimos nem os máximos são ótimos.

\textbf{Degradação em consultas simples: o custo da adaptação agressiva.} Os resultados das consultas simples (Tabela~\ref{tab:simples}) revelam uma limitação importante da adaptação paramétrica: o HyPA-RAG apresenta queda significativa em Article Retrieval Rate (52\%) em comparação ao BM25-only (82\%) e ao Híbrido fixo (70\%). Esta degradação é diretamente atribuível à redução de $top\_k$ para 5 em consultas classificadas como simples, que resulta em subrecuperação de artigos relevantes. A Citation Article Rate também sofre (40\% vs. 50\% do BM25-only e Híbrido fixo). Este resultado evidencia que o classificador heurístico, ao reduzir agressivamente o $top\_k$ para consultas simples, sacrifica cobertura em troca de economia de contexto --- um \textit{tradeoff} que pode não ser desejável quando o custo computacional marginal de fontes adicionais é baixo, como no cenário de execução local com GPU dedicada. Trabalhos futuros devem investigar valores mínimos de $top\_k$ que preservem a cobertura de artigos sem comprometer a eficiência.

\textbf{Tradeoff entre tempo de resposta e adaptabilidade.} O HyPA-RAG apresenta o maior tempo médio de resposta (43,5s), chegando a 96,3s em média nas consultas complexas devido ao $top\_k = 15$ e 3 \textit{rewrites}, contra 26,2s do BM25-only e 23,8s do Híbrido fixo. Este incremento é atribuível ao custo computacional das reformulações de consulta via LLM (para consultas médias e complexas) e ao $k$ elevado em consultas complexas. Entretanto, dado que o cenário de uso primário é a consulta deliberativa por gestores ambientais -- e não a interação conversacional em tempo real --, um tempo de resposta na faixa de 40--50 segundos em média permanece aceitável para o domínio de aplicação.

\subsection{Limitações}

Os resultados apresentados devem ser interpretados à luz de três limitações reconhecidas. Primeiramente, o \textit{dataset} de 100 consultas foi construído especificamente para este estudo e não constitui um \textit{benchmark} padronizado da área, impedindo comparações diretas com outros sistemas. A construção de \textit{benchmarks} públicos para RAG jurídico em português permanece como lacuna relevante na comunidade. Em segundo lugar, a avaliação empregou um único LLM (Mistral 7B), e o desempenho do pipeline pode variar significativamente com modelos de diferentes capacidades (e.g., Llama 3, GPT-4). Terceiro, o classificador de complexidade opera exclusivamente sobre heurísticas linguísticas, sem treinamento supervisionado; um classificador aprendido poderia capturar padrões mais sutis de complexidade. Por fim, o corpus de 6 legislações, embora representativo dos marcos legais ambientais mais consultados, não abrange a totalidade do arcabouço normativo (e.g., decretos regulamentadores, normativas estaduais, portarias do IBAMA).

\section{Conclusão}

Este trabalho apresentou o EcoLex RAG, um sistema HyPA-RAG adaptado ao domínio jurídico-ambiental brasileiro que integra recuperação tripla (BM25 + Legal-BERTimbau/FAISS + Knowledge Graph legal), fusão via Reciprocal Rank Fusion e adaptação paramétrica governada por um classificador de complexidade de consultas. O estudo de ablação com 100 consultas sobre 6 legislações demonstrou que a adaptação paramétrica permite ao sistema manter desempenho competitivo em métricas de recuperação e citação enquanto opera com 32\% menos fontes do que configurações estáticas, evidenciando eficiência superior na gestão da janela de contexto do LLM. O resultado mais expressivo concentrou-se no estrato de consultas de complexidade média, onde o HyPA-RAG alcançou 54\% em Citation Article Rate, superando todas as demais configurações.

As contribuições deste trabalho estendem-se em três direções: (i) a demonstração empírica de que a adaptação paramétrica é uma estratégia viável para otimizar sistemas RAG no domínio legal brasileiro; (ii) a validação do Legal-BERTimbau como modelo de \textit{embedding} adequado para recuperação semântica de trechos legislativos em português; e (iii) a proposta de um pipeline de geração explicável cuja estrutura de citação obrigatória (lei, artigo, parágrafo) alinha o sistema com os requisitos de transparência e auditabilidade da administração pública.

No contexto pós-COP30, em que o Brasil assume compromissos climáticos que demandam governança ambiental ágil e tecnicamente fundamentada, ferramentas como o EcoLex RAG representam uma ponte necessária entre o arcabouço normativo e a prática administrativa. A capacidade de consultar legislação ambiental em linguagem natural e receber respostas explicáveis, com citação exata de dispositivos legais, tem potencial para democratizar o acesso a informação jurídica e qualificar o processo decisório de gestores públicos em todos os níveis federativos.

Como trabalhos futuros, destacam-se: (i) o \textit{fine-tuning} do LLM em corpora jurídico-ambientais brasileiros, visando melhorar a qualidade das citações e a fidelidade ao texto legislativo; (ii) a expansão da base de conhecimento para incluir decretos regulamentadores, normativas estaduais e portarias do IBAMA e ICMBio; (iii) a substituição do classificador heurístico por um modelo aprendido supervisionadamente; (iv) a avaliação com múltiplos LLMs (Llama 3, Sabia, GPT-4) para mensurar o impacto do modelo generativo no desempenho \textit{end-to-end}; e (v) a validação do sistema com gestores ambientais em contexto operacional real, avaliando usabilidade, confiança nas respostas e impacto no tempo de tomada de decisão.

\bibliographystyle{sbc}
\bibliography{references}

\end{document}
